\section{完成情况及性能参数}

\subsection{整体介绍}
作品已形成由开发板、TF 卡和 HDMI 显示器组成的完整多媒体展示系统。基础功能、五类拓展功能与四季动态场景集成在同一固化镜像中，全部功能均在实物上正常运行。镜像写入配置 Flash 后，断电再上电，画面与声音自动启动，无需连接电脑。

\photoplaceholder{系统正面整体展示}{使开发板、TF 卡槽、HDMI 连接和显示器同框；屏幕显示作品的四季夜城或静图，保持画面清晰，避免拍入身份信息。}
\photoplaceholder{系统斜向整体展示}{从约 45 度拍摄设备布局与线缆连接，保证开发板与显示内容同时可辨，不遮挡板载按键。}

\subsection{工程成果}

\subsubsection{硬件成果}
系统使用 HX4S20C 板载资源构成 FPGA、TF、SDRAM、配置 Flash 和 HDMI 接口。图片和音乐从 TF 卡进入 FPGA，图像经 SDRAM 缓存后与动态场景及各图层合成，经 HDMI 同时送出画面与声音。全部运行时处理在 FPGA 内完成，无需外接处理器。

\photoplaceholder{开发板与接口实物}{近距离拍摄开发板正面，使芯片、TF 卡槽、HDMI 输出、按键与拨码开关的位置清楚可见。}

\subsubsection{FPGA 逻辑成果}
FPGA 逻辑涵盖 TF 卡访问与素材解析、双线性缩放、六帧缓存与刷新调度、音频播放与音量处理、图层合成、中文菜单与字幕、四季场景生成和 720p 音视频输出，在 EG4S20 上完成综合、布局布线并全部时序收敛。

\photoplaceholder{工程工具与构建结果界面}{在 TD 工程工具中显示当前作品的器件、主要资源或布局布线结果，截图时保留可辨读的结果区域，裁去本机用户名、路径和与作品无关的信息。}

\begin{center}\small
\begin{tabular}{p{4.6cm}p{2.4cm}p{2.4cm}p{3.2cm}}
\toprule
资源 & 使用 & 器件容量 & 利用率\\\midrule
LUT & 16454 & 19600 & 84.0\%\\
寄存器 & 9007 & 19600 & 46.0\%\\
SLICE & 8533 & 9800 & 87.1\%\\
9 Kb RAM 块 & 63 & 64 & 98.4\%\\
32 Kb RAM 块 & 16 & 16 & 100\%\\
DSP & 24 & 29 & 82.8\%\\
PLL & 4 & 4 & 100\%\\
\bottomrule
\end{tabular}
\end{center}
片上 RAM 与 PLL 全部用满、LUT 利用率达 84\%，说明作品充分挖掘了 EG4S20 的硬件资源；在如此高的利用率下全部时钟域仍然时序收敛，setup 最差余量 $+0.682$ ns，保证了系统长期稳定运行。

\subsection{特性成果}
\subsubsection{基础功能与扩展功能总览}
\begin{center}\small
\begin{tabular}{p{2.8cm}>{\raggedright\arraybackslash}p{7.4cm}p{3.0cm}}
\toprule
功能 & 作品实现 & 完成情况\\\midrule
图片播放 & TF 卡读取 24 位 BMP，经 SDRAM 缓存后 HDMI 720p 输出 & 已完成，实物运行正常\\
手动与自动切图 & 按键前后切换；拨码开启自动轮播，周期 \tablegroup{3 / 5 / 15} 秒可选 & 已完成，实物运行正常\\
音频播放 & 48 kHz 立体声 PCM，HDMI 音频输出，四首曲目菜单选择 & 已完成，实物运行正常\\
固化启动 & 统一镜像写入 Flash，上电自动加载运行 & 已完成，实物运行正常\\
文字叠加 & 七条中文字幕随\tablegroup{图片 / 场景}切换并滚动显示 & 已完成，实物运行正常\\
图像转场 & 旧画面淡出至黑、新画面淡入 & 已完成，实物运行正常\\
不同尺寸缩放 & BMP 几何解析、等比双线性缩放与居中补黑 & 已完成，实物运行正常\\
实时参数调节 & 中文菜单调节亮度、对比度、音量、字幕、轮播等 & 已完成，实物运行正常\\
音频可视化 & 真实 PCM 驱动主画面 64 柱与街区 12 柱音乐屏 & 已完成，实物运行正常\\
四季动态场景 & 四景切换、天气灯光实时变化、无缝返回静图 & 已完成，实物运行正常\\
\bottomrule
\end{tabular}
\end{center}

\subsubsection{图片读取、缩放与轮播}
系统配置六幅展示图，包括腕表、饮品、芯片、暴雨提示、车道提示及课堂展示。其中芯片图为 $640\times480$、24 位 BMP，系统自动等比缩放到 $480\times360$ 内容区并在两侧补黑，原图比例完整保留；其他尺寸的图片由同一套几何规则自动适配。按键切图与自动轮播都伴随淡出淡入转场，背景音乐持续播放。

\photoplaceholder{图片显示与比例适配实物}{拍摄芯片图完整画面及其左右黑边；另拍一幅宽画面素材，能看出显示区域与原始长宽比的关系。}
\photoplaceholder{手动切图与自动轮播实物}{分别拍摄操作按键／拨码与切换前后的画面，若需说明转场过程，可从实拍视频中选取淡出、黑场和淡入三帧。}

\subsubsection{中文菜单与实时参数}
拨动 SW2 打开中文菜单，方向键选择功能，确认键进入调整；调整结果立即作用于画面或声音。亮度和对比度改变背景像素，音量同时改变声音和音频柱高度，观众可直接观察控制量与效果的对应关系。

\photoplaceholder{中文菜单与参数调节实物}{拍摄一级菜单、一个二级参数页及调整前后画面，保证参数文字和背景可同时辨认；必要时分别放置两张照片。}

\subsubsection{滚动字幕与图像转场}
每幅静图和四季场景都有对应的中文字幕，切换内容时字幕同步更新并重新滚动，可在菜单中开关。切图时旧画面平滑淡出到黑，再淡入新图，过渡自然流畅。

\photoplaceholder{中文字幕与转场实物}{拍摄字幕清晰可读的两幅不同画面；转场可用连续实拍帧展示旧图、黑场与新图，勿用后期效果替代设备输出。}

\subsubsection{音频播放与音频可视化}
系统播放 48 kHz 立体声音乐，可在菜单中选择四首曲目。主画面底部 64 根幅度柱随音乐节奏起伏滚动，夜城街区的音乐屏同步跳动；调节音量时，声音与柱高一同变化。

\photoplaceholder{选曲与音频可视化实物}{画面同时包含曲目菜单或播放画面与音频柱；可另外拍摄静音和恢复音量两个状态，并通过实录声音保存听音证据。}

\subsubsection{四季动态夜城}
拨动 SW3 进入动态场景，方向键在雨夜霓虹、黄昏雷云、薄雾深夜和飘雪清晨之间切换。同一座城市在四季中呈现截然不同的色彩与天气：雨线、闪电、雾带、雪粒持续运动，霓虹招牌呼吸闪烁，街区音乐屏随音乐跳动。再次拨动 SW3 即回到原静态图片，音乐全程不中断。

\photoplaceholder{四季动态场景实物组图}{以相同机位分别拍摄四种场景；每张图保留清晰的屏幕边界，画面能区分雨、云、雾和雪的视觉变化。}
\photoplaceholder{街区细节与音乐联动实物}{近拍 NIGHT CITY 广告屏、店铺招牌、灯带及音频小屏；以视频或不同时间的实拍帧记录动态变化。}

\subsubsection{仿真验证结果}
四季场景的 28 帧 RTL 输出与独立参考模型逐像素一致，天气变化和音频柱历史符合设计；字幕验证覆盖全部七条文案和 214 个字形；$640\times480$ 输入的缩放结果在全部 230400 个画布像素上与模型的最大通道偏差仅为 1。仿真、布局布线和实物运行三级验证共同确认了设计的正确性。

\subsubsection{固化与上电运行}
最终镜像写入配置 Flash 并通过自动校验。断电后重新上电，无需重新下载或按复位键，系统自动完成素材加载，画面与声音正常启动，全部功能均可在该镜像上直接使用。

\photoplaceholder{固化上电后的运行实物}{拍摄上电运行的开发板与显示器；保留不连接下载电脑的正常使用画面，条件允许时用连续视频补充断电、上电与启动过程。}
